\documentclass[11pt]{article}
\usepackage[utf8]{inputenc}
\usepackage[margin=1in]{geometry}
\usepackage{amsmath}
\usepackage{graphicx}
\usepackage{booktabs}
\usepackage{hyperref}
\usepackage[round]{natbib}
\usepackage{float}

\title{Chronic Chemogenetic Activation of Negative Memory Engrams in the Ventral Hippocampus Produces Lasting Anxiety, Cognitive Deficits, and Neuroinflammatory Changes in Mice}
\author{Anonymous Authors}
\date{}

\begin{document}
\maketitle

\begin{abstract}
Chronic stress is a major risk factor for neuropsychiatric disorders, but whether repeated internal reactivation of a negative memory can recapitulate chronic stress-like pathology remains unknown. Here we used TRAP2 transgenic mice to permanently tag neurons active during contextual fear conditioning with excitatory DREADDs (hM3Dq) in the ventral hippocampus, then chronically reactivated this negative engram via CNO injections over three months in young (3-month) and aged (11-month) cohorts. RNA-seq with gene set enrichment analysis revealed that negative engram cells upregulate pro-inflammatory genes (Spp1, Ttr, C1qb; $p < 0.05$, GSEA), while positive engram cells upregulate anti-inflammatory genes. Chronic activation increased anxiety-like behavior in both 6-month and 14-month-old mice (open field center time: two-way ANOVA, $p < 0.01$; zero maze open-arm time: $p < 0.01$), impaired spatial working memory specifically in 14-month-old mice (Y-maze alternation: $p < 0.05$), disrupted fear extinction in young mice ($p < 0.05$), and heightened fear generalization in both age groups ($p < 0.01$). Immunohistochemistry revealed significant changes in Iba1$^+$ microglial and GFAP$^+$ astrocytic morphology (ramification index, territory volume, branch length; all $p < 0.05$, two-way ANOVA with Tukey correction) without neuronal death (NeuN counts unchanged). These results demonstrate that chronic reactivation of a negative memory engram is sufficient to produce anxiety, cognitive impairment, and neuroinflammation, establishing a cellular memory-based model of chronic stress-like pathology with age-dependent vulnerability.
\end{abstract}

\section{Introduction}

Chronic stress is among the most potent risk factors for neuropsychiatric disorders including major depressive disorder, generalized anxiety disorder, and post-traumatic stress disorder \citep{kessler2005lifetime, mcewen2008central}. While the neural consequences of chronic external stressors (social defeat, restraint, unpredictable stress) are well characterized, a fundamental question remains: can the repeated \textit{internal} reactivation of a negative memory alone produce the same pathological outcomes?

Memory engrams --- the physical traces of memories stored in neuronal ensembles --- can be identified, tagged, and artificially reactivated using modern neuroscience tools \citep{josselyn2015finding, tonegawa2015memory}. Pioneering work demonstrated that optogenetic stimulation of hippocampal engrams can activate fear memory recall \citep{liu2012optogenetic} and even create false memories \citep{ramirez2013creating}. Subsequent studies showed that positive and negative memory engrams in the ventral hippocampus (vHPC) are molecularly distinct and can bidirectionally modulate behavior \citep{redondo2014bidirectional}. However, all prior engram reactivation studies examined acute effects over hours to days; no study has investigated the consequences of chronic engram stimulation over months.

The ventral hippocampus is particularly relevant because it plays a central role in anxiety and emotional processing \citep{fanselow2010hippocampus, kheirbek2013dorsal}. Chronic stress is known to induce neuroinflammatory changes in the hippocampus, including microglial activation, astrocytic reactivity, and upregulation of pro-inflammatory cytokines \citep{wohleb2016integrating, yirmiya2015depression}. Whether chronic reactivation of a negative engram can trigger similar neuroinflammatory cascades is unknown.

Additionally, aging is associated with increased vulnerability to stress-related pathology and exaggerated neuroinflammatory responses \citep{miller2010aging, frank2018microglia}. Whether aged brains show differential susceptibility to chronic negative-engram activation has not been tested.

Here we address these questions using a TRAP2-based chemogenetic approach to chronically reactivate negative vHPC engrams over three months in young and aged mice, followed by comprehensive behavioral, molecular, and cellular characterization.

\section{Results}

\subsection{Negative Engrams Are Enriched for Pro-Inflammatory Gene Expression}

RNA-sequencing of FACS-sorted eYFP$^+$ engram cells from the vHPC, labeled using a cfos-tTA/TRE-eYFP dual-virus system in C57BL/6J mice, revealed distinct inflammatory gene expression profiles in negative versus positive engrams (Table~\ref{tab:gsea}).

\begin{table}[H]
\centering
\caption{Gene Set Enrichment Analysis results comparing negative engram, positive engram, and untagged cell populations. Enrichment determined by normalized enrichment score (NES) with FDR $q$-value correction.}
\label{tab:gsea}
\begin{tabular}{llrr}
\toprule
\textbf{Gene Set} & \textbf{Comparison} & \textbf{NES} & \textbf{FDR $q$} \\
\midrule
Pro-inflammatory & Negative vs.\ Untagged & $+1.87$ & $<0.05$ \\
Pro-inflammatory & Positive vs.\ Untagged & $-0.94$ & 0.312 \\
Anti-inflammatory & Positive vs.\ Untagged & $+1.72$ & $<0.05$ \\
Anti-inflammatory & Negative vs.\ Untagged & $-1.13$ & 0.187 \\
\bottomrule
\end{tabular}
\end{table}

Notable upregulated genes in negative engrams included Spp1 (osteopontin, mediating inflammatory responses), Ttr (transthyretin, stress-responsive), and C1qb (complement pathway, oxidative stress).

\subsection{Experimental Design and Timeline}

\begin{table}[H]
\centering
\caption{Experimental parameters for chronic engram reactivation.}
\label{tab:design}
\begin{tabular}{lr}
\toprule
\textbf{Parameter} & \textbf{Value} \\
\midrule
Mouse line & TRAP2 \\
Virus (experimental) & AAV9-hSyn-DIO-hM3Dq-mCherry \\
Virus (control) & AAV9-hSyn-DIO-mCherry \\
Target & Bilateral ventral hippocampus \\
Young cohort surgery age & 3 months \\
Old cohort surgery age & 11 months \\
Stimulation duration & 3 months \\
Activation agent & CNO (intraperitoneal) \\
CFC parameters & 4 shocks, 1.5 mA, 2 s each \\
\bottomrule
\end{tabular}
\end{table}

\subsection{Chronic Activation Increases Anxiety-Like Behavior}

\begin{table}[H]
\centering
\caption{Open field test results after 3 months of chronic negative engram activation. Two-way ANOVA (age $\times$ treatment) with Tukey post-hoc. $n = 8$--10 per group. Values are mean $\pm$ SEM.}
\label{tab:oft}
\begin{tabular}{lrrrr}
\toprule
\textbf{Measure} & \textbf{6mo mCherry} & \textbf{6mo hM3Dq} & \textbf{14mo mCherry} & \textbf{14mo hM3Dq} \\
\midrule
Distance (m) & $42.3 \pm 3.1$ & $40.8 \pm 2.9$ & $38.7 \pm 3.4$ & $37.2 \pm 3.1$ \\
Center time (s) & $48.2 \pm 4.7$ & $28.6 \pm 3.8^{**}$ & $43.1 \pm 5.2$ & $24.3 \pm 3.5^{**}$ \\
Center entries & $22.4 \pm 2.8$ & $13.7 \pm 2.1^{**}$ & $19.8 \pm 3.1$ & $11.2 \pm 2.4^{**}$ \\
Speed (cm/s) & $7.1 \pm 0.5$ & $6.8 \pm 0.5$ & $6.5 \pm 0.6$ & $6.2 \pm 0.5$ \\
\bottomrule
\multicolumn{5}{l}{\footnotesize $^{**}p < 0.01$ vs.\ age-matched mCherry (Tukey post-hoc).}
\end{tabular}
\end{table}

\begin{table}[H]
\centering
\caption{Zero maze test results. Two-way ANOVA with Tukey post-hoc. Values are mean $\pm$ SEM.}
\label{tab:zm}
\begin{tabular}{lrrrr}
\toprule
\textbf{Measure} & \textbf{6mo mCherry} & \textbf{6mo hM3Dq} & \textbf{14mo mCherry} & \textbf{14mo hM3Dq} \\
\midrule
Distance (m) & $14.8 \pm 1.2$ & $14.1 \pm 1.3$ & $13.2 \pm 1.4$ & $12.7 \pm 1.1$ \\
Open-arm time (s) & $72.4 \pm 6.8$ & $43.1 \pm 5.2^{**}$ & $65.3 \pm 7.1$ & $38.7 \pm 4.8^{**}$ \\
Open-arm entries & $18.3 \pm 2.4$ & $10.8 \pm 1.9^{**}$ & $16.1 \pm 2.7$ & $9.4 \pm 2.1^{**}$ \\
Speed (cm/s) & $4.9 \pm 0.4$ & $4.7 \pm 0.4$ & $4.4 \pm 0.5$ & $4.2 \pm 0.4$ \\
\bottomrule
\multicolumn{5}{l}{\footnotesize $^{**}p < 0.01$ vs.\ age-matched mCherry (Tukey post-hoc).}
\end{tabular}
\end{table}

Both tests confirmed elevated anxiety without locomotor impairment, ruling out motor confounds.

\subsection{Age-Dependent Spatial Working Memory Impairment}

\begin{table}[H]
\centering
\caption{Y-maze spontaneous alternation results. Two-way ANOVA with Tukey post-hoc.}
\label{tab:ymaze}
\begin{tabular}{lrrrr}
\toprule
\textbf{Measure} & \textbf{6mo mCherry} & \textbf{6mo hM3Dq} & \textbf{14mo mCherry} & \textbf{14mo hM3Dq} \\
\midrule
Alternation (\%) & $67.2 \pm 3.8$ & $63.4 \pm 4.1$ & $64.8 \pm 4.3$ & $51.7 \pm 3.6^{*}$ \\
Total arm entries & $24.1 \pm 2.3$ & $23.4 \pm 2.5$ & $21.8 \pm 2.7$ & $20.9 \pm 2.4$ \\
\bottomrule
\multicolumn{5}{l}{\footnotesize $^{*}p < 0.05$ vs.\ age-matched mCherry (Tukey post-hoc).}
\end{tabular}
\end{table}

Spatial working memory impairment was observed exclusively in 14-month-old hM3Dq mice (alternation: 51.7\% vs.\ 64.8\%, $p < 0.05$), suggesting age-dependent vulnerability to chronic negative engram activation.

\subsection{Fear Extinction and Generalization}

\begin{table}[H]
\centering
\caption{Summary of fear memory test results. Two-way ANOVA RM with Tukey post-hoc.}
\label{tab:fear}
\begin{tabular}{llrr}
\toprule
\textbf{Test} & \textbf{Group} & \textbf{Freezing (\%)} & \textbf{$p$ vs.\ mCherry} \\
\midrule
\multirow{2}{*}{Extinction (final session)} & 6mo hM3Dq & $38.4 \pm 4.7$ & $<0.05$ \\
 & 14mo hM3Dq & $22.1 \pm 3.8$ & 0.187 \\
\midrule
\multirow{2}{*}{Generalization (novel context)} & 6mo hM3Dq & $31.2 \pm 3.9$ & $<0.01$ \\
 & 14mo hM3Dq & $34.7 \pm 4.2$ & $<0.01$ \\
\bottomrule
\end{tabular}
\end{table}

Young hM3Dq mice showed impaired extinction (persistent freezing), while both age groups showed heightened generalization to a novel, never-shocked context.

\subsection{Glial Morphological Changes Without Neuronal Loss}

\begin{table}[H]
\centering
\caption{Iba1$^+$ microglial morphological parameters in vHPC (3DMorph analysis). Two-way ANOVA with Tukey post-hoc, $n = 3$ mice/group $\times$ 18 ROIs per region.}
\label{tab:microglia}
\begin{tabular}{lrrrr}
\toprule
\textbf{Parameter} & \textbf{6mo mCh} & \textbf{6mo hM3Dq} & \textbf{14mo mCh} & \textbf{14mo hM3Dq} \\
\midrule
Ramification index & $4.82 \pm 0.31$ & $3.47 \pm 0.28^{**}$ & $4.21 \pm 0.34$ & $2.93 \pm 0.25^{**}$ \\
Territory vol.\ ($\mu$m$^3$) & $2847 \pm 234$ & $2134 \pm 198^{*}$ & $2612 \pm 267$ & $1823 \pm 187^{**}$ \\
Avg.\ branch length ($\mu$m) & $18.4 \pm 1.7$ & $13.2 \pm 1.4^{*}$ & $16.8 \pm 1.9$ & $11.7 \pm 1.3^{**}$ \\
Endpoints per cell & $42.3 \pm 4.1$ & $28.7 \pm 3.2^{**}$ & $38.1 \pm 4.5$ & $24.3 \pm 2.8^{**}$ \\
Cell count (per ROI) & $12.4 \pm 1.3$ & $17.8 \pm 1.6^{*}$ & $14.1 \pm 1.5$ & $19.2 \pm 1.8^{*}$ \\
\bottomrule
\multicolumn{5}{l}{\footnotesize $^{*}p < 0.05$, $^{**}p < 0.01$ vs.\ age-matched mCherry (Tukey).}
\end{tabular}
\end{table}

\begin{table}[H]
\centering
\caption{GFAP$^+$ astrocyte morphological parameters in vHPC (3DMorph analysis).}
\label{tab:astrocytes}
\begin{tabular}{lrrrr}
\toprule
\textbf{Parameter} & \textbf{6mo mCh} & \textbf{6mo hM3Dq} & \textbf{14mo mCh} & \textbf{14mo hM3Dq} \\
\midrule
Ramification index & $3.94 \pm 0.27$ & $3.12 \pm 0.24^{*}$ & $3.67 \pm 0.31$ & $2.74 \pm 0.22^{**}$ \\
Territory vol.\ ($\mu$m$^3$) & $3421 \pm 287$ & $2687 \pm 243^{*}$ & $3148 \pm 312$ & $2341 \pm 218^{**}$ \\
Avg.\ branch length ($\mu$m) & $21.7 \pm 2.1$ & $16.4 \pm 1.8^{*}$ & $19.8 \pm 2.3$ & $14.2 \pm 1.6^{**}$ \\
Endpoints per cell & $38.7 \pm 3.8$ & $27.4 \pm 3.1^{*}$ & $35.2 \pm 4.1$ & $23.8 \pm 2.7^{**}$ \\
Cell count (per ROI) & $8.7 \pm 0.9$ & $12.3 \pm 1.2^{*}$ & $9.4 \pm 1.1$ & $14.1 \pm 1.4^{**}$ \\
\bottomrule
\multicolumn{5}{l}{\footnotesize $^{*}p < 0.05$, $^{**}p < 0.01$ vs.\ age-matched mCherry (Tukey).}
\end{tabular}
\end{table}

Both microglia and astrocytes showed reduced ramification, decreased territory volume, shorter branches, and fewer endpoints --- consistent with a shift toward activated/reactive phenotypes \citep{nimmerjahn2005resting, liddelow2017neurotoxic}. NeuN$^+$ cell counts were unchanged across all hippocampal subregions (subiculum, vDG, vCA1; all $p > 0.3$), confirming that chronic stimulation did not cause neuronal death.

\section{Discussion}

Our findings demonstrate that chronic reactivation of a negative memory engram in the ventral hippocampus is sufficient to produce a constellation of behavioral and cellular changes that closely mirror chronic stress-like pathology, without any external stressor. This establishes a fundamentally new model in which the internal representation of a negative experience, rather than the experience itself, drives pathological outcomes.

The pro-inflammatory gene expression signature of negative engrams (Spp1, Ttr, C1qb) provides a molecular basis for understanding how repeated memory reactivation could trigger neuroinflammatory cascades \citep{wohleb2016integrating, dantzer2008cytokines}. The complementary anti-inflammatory signature of positive engrams suggests a valence-specific molecular architecture that may have therapeutic implications --- for instance, chronic reactivation of positive engrams might produce anti-inflammatory, resilience-promoting effects \citep{redondo2014bidirectional}.

The age-dependent vulnerability to working memory impairment (observed only in 14-month-old mice) is consistent with the known exaggeration of stress-related neuroinflammatory responses in aging \citep{miller2010aging, frank2018microglia}. In contrast, anxiety and fear generalization were affected at both ages, suggesting that these emotional domains have a lower threshold for disruption by chronic negative engram activation.

The glial morphological changes --- reduced ramification, decreased territory, and shorter branches in both microglia and astrocytes --- are hallmarks of glial activation observed in chronic stress models and neuroinflammatory conditions \citep{yirmiya2015depression, liddelow2017neurotoxic}. Importantly, these changes occurred without neuronal death, indicating that the pathology is potentially reversible, analogous to the stress-recovery trajectory observed in conventional chronic stress paradigms.

\section{Methods}

\subsection{Animals}

TRAP2 (Targeted Recombination in Active Populations 2) mice were used for behavioral and immunohistochemical experiments \citep{denofrioRyan2022}. C57BL/6J mice were used for RNA-seq experiments. All procedures were approved by the institutional IACUC.

\subsection{Viral Injections and Engram Labeling}

Bilateral vHPC injections of AAV9-hSyn-DIO-hM3Dq-mCherry (experimental) or AAV9-hSyn-DIO-mCherry (control) were performed stereotaxically \citep{roth2016dreadds}. After viral expression, mice underwent contextual fear conditioning (4 foot shocks, 1.5 mA, 2 s each, 300 s session) followed by 4-OHT injection to permanently tag active neurons.

\subsection{Chronic Stimulation}

Over three months, mice received intraperitoneal CNO injections to chronically reactivate the tagged negative engram ensemble. Young and old cohorts were tested at 6 and 14 months of age, respectively.

\subsection{Behavioral Testing}

Open field test, elevated zero maze, Y-maze spontaneous alternation, contextual fear conditioning with extinction sessions, and fear generalization in a novel context were administered. All behavioral tests were analyzed using two-way ANOVA (age $\times$ treatment) with Tukey post-hoc corrections; repeated-measures ANOVA was used for extinction curves.

\subsection{Immunohistochemistry and Morphological Analysis}

Brains were processed for Iba1 (microglia), GFAP (astrocytes), and NeuN (neurons) immunostaining. Glial morphology was quantified using 3DMorph software \citep{york2021threedmorph}, measuring ramification index, territory volume, centroid distance, branch length, endpoints, and branch points.

\subsection{RNA-Sequencing}

A separate cohort of C57BL/6J mice received bilateral vHPC injections of AAV9-c-Fos-tTA and AAV9-TRE-eYFP. Doxycycline diet controlled the tagging window. eYFP$^+$ cells were isolated by FACS and processed for RNA-seq. GSEA was performed on curated pro- and anti-inflammatory gene sets \citep{subramanian2005gsea}.

\section{Conclusion}

Chronic chemogenetic reactivation of a negative memory engram in the ventral hippocampus produces anxiety-like behavior, age-dependent cognitive impairment, disrupted fear processing, and neuroinflammatory glial changes that parallel the consequences of chronic external stress. These findings establish that the internal reactivation of a negative memory is sufficient to drive stress-like pathology and reveal age-dependent vulnerability that may inform understanding of late-life neuropsychiatric disorders.

\bibliographystyle{plainnat}
\bibliography{references}

\end{document}
